
\chapter{Open Questions and Emerging Directions}
\section{Current Research Directions}
\begin{itemize}
    \item Hardware/software hybrid TM.
    \item Non-blocking algorithms vs. TM.
    \item Persistent memory TM.
    \item Distributed TM.
    \item GPU TM (see Chapter~\ref{chap:gpu-tm}): no vendor hardware,
          software conflict detection, and the search for scalable commit
          protocols at thousands of concurrent transactions.
    \item Persistent TM: durable commit with volatile metadata; crash-consistency without a database's write-ahead log.
    \item Distributed TM: exporting the TM API over replicated or geo-distributed storage (the companion repo wraps a TiKV cluster).
    \item Hybrid systems: HTM fast path, STM (or serialized) fallback, and the theory of `phased' and `shipwreck' hybrids.
    \item Verified/certified STMs: mechanically checked backends, and the growing standard of TLA+-verified designs.
    \item Energy-aware TM: correlation between abort rate and power on mobile/embedded cores.
\end{itemize}

The implementation boundary is now part of the research question. The companion
repo has working CPU backends from SGL through SwissTM, TinySTM variants,
NOrec-BF, TL2, TSC-TM, MVLog, persistent and distributed SGL, and a TiKV
Percolator-style backend; it also contains GPU designs such as PR-STM, CSMV,
GPUTX, GAccO, GPU-JVSTM, and GUST. All are tested against the same transactional
and data-structure suites, and the core protocols have TLA+ models. This makes
the open questions unusually concrete: the same bank invariant should survive
hardware aborts, validation filters, crash recovery, remote prepare/commit, and
GPU commit logs.

\begin{table}[htbp]
\centering
\begin{tabular}{p{0.18\textwidth}p{0.25\textwidth}p{0.30\textwidth}p{0.14\textwidth}}
\hline
Direction & Representative backends & Main pressure point & Repo signal \\
\hline
Hybrid HTM/STM & TSX-SGL, phased HTM-to-STM sketches & Capacity aborts, ownership agreement between paths, fallback cliffs & Tested; no portable HTM standard \\
Classic STM & NOrec, NOrec-BF, TL2, TSC-TM, TinySTM, SwissTM & Validation cost, metadata contention, memory reclamation & Bank fuzz: TL2 about 2.06M txns/s, TSC-TM about 2.05M on M1 Pro \\
Persistent TM & NV-HTM, Romulus, PersistentSGL & Durable ordering, redo/undo idempotence, volatile metadata after crash & Crash invariants modelled \\
Distributed TM & DistributedSGL, TiKV Percolator-2PC wrapper & Remote validation, partial failure, coordinator recovery & Same API, higher-latency commit path \\
GPU TM & PR-STM, CSMV, GPUTX, GAccO, GPU-JVSTM, GUST & No cache-coherent TM hardware, thousands of concurrent transactions & GPU results are device-dependent \\
Multi-version TM & JVSTM, MVLog, GPU-JVSTM & Version retention, read-mostly scalability, reclamation & MVLog about 275K write-heavy and 109K read-mostly txns/s on M1 Pro \\
\hline
\end{tabular}
\caption{Emerging TM directions as implementation pressure points. Numbers are the companion repo's fuzz/bank measurements where CPU measurements are available.}
\label{tab:emerging-directions}
\end{table}

\section{Challenges}
\begin{itemize}
    \item Fairness, progress guarantees (starvation freedom, obstruction freedom).
    \item Programming model adoption.
    \item Legacy code integration.
    \item Performance predictability: HTM's capacity aborts and STM's validation costs make best/worst-case analysis hard.
    \item The missing standard: no widely adopted language integration after the C++ TS; education must bridge this gap.
    \item Tooling: debuggers, profilers, and race detectors are lock-oriented; TM tools are nascent by comparison.
\end{itemize}

\section{Worked Example: A Hybrid, In Miniature}
\index{hybrid TM}\index{phased hybrid}\index{fallback}\index{capacity abort}\index{wasted work}
\label{sec:hybrid-example}

The ``hardware fast path, software fallback'' idea, small enough to see whole:

\begin{lstlisting}[language=C, caption={A minimal phased hybrid: try RTM, fall back to the STM.}]
void tx_region(void) {
    for (int i = 0; i < MAX_RTM_TRIES; ++i) {
        unsigned s = _xbegin();
        if (s == _XBEGIN_STARTED) {
            body(); _xend(); return;         /* fast path  */
        }
        if (s == _XABORT_CAPACITY) break;    /* stop wasting */
    }
    stm_begin(); body(); stm_commit();       /* slow path  */
}
\end{lstlisting}

\begin{itemize}
    \item This is \S\ref{sec:rtm-example} plus one idea: the fallback is a full STM instead of a global lock, so concurrent commits survive the fallback (no serialisation cliff).
    \item The capacity check turns a wasted-retry storm into a clean handover --- the performance chapter's ``wasted work'' metric in miniature.
    \item The ``phased'' theory of the bullets above governs exactly when to flip; the shipwreck hybrids time the flip adversarially.
    \item Exercise~1 asks whether this is more livelock-prone than pure STM --- the answer turns on whether the hardware and software paths agree on conflict ownership.
\end{itemize}

\begin{figure}[htbp]
\centering
\resizebox{\textwidth}{!}{%
\begin{tikzpicture}[
    node distance=15mm and 23mm,
    >=Stealth,
    state/.style={draw, rounded corners, align=center, minimum width=31mm, minimum height=9mm}
]
\node[state] (start) {transaction\\request};
\node[state, right=of start] (rtm) {RTM attempt};
\node[state, right=of rtm] (hwcommit) {hardware\\commit};
\node[state, below=of rtm] (classify) {abort\\classification};
\node[state, left=of classify] (retry) {retry budget\\remaining};
\node[state, right=of classify] (stm) {STM fallback};
\node[state, right=of stm] (swcommit) {software\\commit};

\draw[->] (start) -- (rtm);
\draw[->] (rtm) -- node[above]{success} (hwcommit);
\draw[->] (rtm) -- node[right]{abort} (classify);
\draw[->] (classify) -- node[above]{transient} (retry);
\draw[->] (retry) |- (rtm);
\draw[->] (classify) -- node[above]{capacity or budget exhausted} (stm);
\draw[->] (stm) -- (swcommit);
\end{tikzpicture}
}%

\caption{A phased hybrid state machine. The research problem is not the branch itself, but the ownership contract between the RTM path and the STM fallback.}
\label{fig:phased-hybrid-state}
\end{figure}
\section{Summary}
\begin{itemize}
    \item TM is a rich research field with many unsolved problems.
    \item The student should leave with a critical understanding of where TM stands and where it might go.
    \item Directions worth watching divide into three families: hardware (hybrid HTM, GPU), systems (persistent, distributed), and methods (verification, education).
    \item Read this chapter as a \emph{synthesis organised around where the book's own codebase stops working}, not as a disconnected wish-list. The limits that surface in each family are the same limits the earlier parts hit: HTM's capacity aborts (Ch.~16), STM's validation and reclamation costs (Ch.~9, 13), the missing compiler/language contract (Ch.~11), and the absence of coherence on GPUs (Ch.~17). Each open direction in this chapter is one of those known limits restated as a research target.
\end{itemize}

\section{Exercises}
\begin{enumerate}
    \item Evaluate whether a hybrid TM that starts in hardware and falls back to software is inherently more susceptible to livelock than a pure-STM system.
    \item Propose a minimal extension to a persistent-memory TM that ensures durability of committed transactions across power failures.
    \item Argue for or against: “Transactional memory will eventually replace locks as the default concurrency primitive in systems programming.”
    \item \emph{[theory]} Classify the phased hybrid in Figure~\ref{fig:phased-hybrid-state} under obstruction freedom, starvation freedom, and lock freedom. State which assumptions about retry bounds, fallback admission, and conflict ownership are necessary for each classification.
\end{enumerate}

%  APPENDICES
